\documentclass[11pt]{article}
\usepackage[margin=1in]{geometry}
\usepackage{amsmath,amssymb,amsthm}
\usepackage[T1]{fontenc}
\usepackage{lmodern}
\usepackage[hidelinks]{hyperref}
\usepackage{xurl}
\emergencystretch=2em
\newtheorem{theorem}{Theorem}[section]
\newtheorem{proposition}[theorem]{Proposition}
\newtheorem{corollary}[theorem]{Corollary}
\theoremstyle{remark}\newtheorem{remark}[theorem]{Remark}
\newcommand{\Ric}{\operatorname{Ric}}
\newcommand{\Rm}{\operatorname{Rm}}
\newcommand{\R}{\mathbb R}
\newcommand{\ip}[2]{\langle #1,#2\rangle}
\newcommand{\leanref}[2]{\par\smallskip\noindent{\small\textit{Lean declaration.}\ \path{#1}\par\noindent\textit{Module.}\ \path{#2}}\par\smallskip}
\title{A Lean Formalization of Uniform Curve Estimates under Ricci Flow\\[0.5em]
\large A module for the Perelman--Morgan--Tian finite extinction argument}
\author{GPT-6 Astra \and Juii-hang Leung}
\date{September 9, 2026}
\begin{document}
\maketitle
\begin{abstract}
We describe a Lean formalization of length and total-curvature estimates for
smooth immersed closed curves evolving by curve shortening in an ambient Ricci
flow. The development derives the corrected curvature evolution from actual
geometric definitions, treats curvature zeros by regularization, and integrates
against the moving arclength measure. For a compact base Ricci flow it constructs
finite ambient tensor bounds and proves that the same coefficients apply to
products with every supplied static one-dimensional Riemannian metric. The
constants are chosen before the factor metrics, curve families, and initial
bounds. The result is an interface for the finite extinction argument, conditional
on existence of the smooth flows and on uniform initial data. It does not prove
flow existence or finite extinction. We explain the relation to the 2015
Morgan--Tian correction, the use of a spatial time-dependent connection, the
precise library interfaces, and the scope of the checked artifact.
\end{abstract}

\section{Scope and mathematical setting}
Perelman's finite extinction argument uses controlled curve families in an
ambient Ricci flow \cite{perelman2003}. The relevant estimates were developed in
Chapter 19 of Morgan and Tian's account \cite{morgantian2007}; their correction
\cite{morgantian2015} changes the curvature evolution and the dependence on initial
data. The formalization reported here concerns this estimate module. It supplies
neither the surrounding topological argument nor the analytic existence theory.

All metrics and connections below are genuine smooth Riemannian objects. Curves
are represented by maps $c(t,x)$ with $x\in\R$ and period one; integration over
$[0,1]$ therefore represents integration over the parameter circle. An immersion
means $c_x(t,x)\ne0$. Write
\[
 X=c_x,\qquad v=|X|_{g(t)},\qquad S=v^{-1}X,\qquad
 H=D_sS,\qquad q=|H|_{g(t)}^2,\qquad k=\sqrt q,
\]
where $D_s=v^{-1}D_x$ is the spatial Levi--Civita derivative along the curve. The
curve-shortening equation is $c_t=H$, while the ambient equation is
$\partial_tg=-2\Ric_g$. The geometric quantities controlled are
\[
 L(t)=\int_0^1v(t,x)\,dx,\quad
 \Theta(t)=\int_0^1 k(t,x)v(t,x)\,dx,\quad
 E(t)=\int_0^1q(t,x)v(t,x)\,dx.
\]
These definitions use the evolving metric, rather than the fixed norm of a chart
model. No global chart containing the curve is assumed; self-intersections are
allowed.

The finite-dimensional manifolds are smooth, boundaryless, Hausdorff, and
sigma-compact; their models are real inner-product spaces of positive dimension.
The principal ambient manifold is compact. Precise Lean typeclass assumptions,
including the time-set convention, are recorded in Appendix~\ref{app:interfaces}.
The metrics are jointly smooth on the specified closed time interval in the
library's within-smoothness sense. The supplied curve flows are jointly smooth,
closed, immersed, and satisfy the actual within-derivative flow equation. These
conditions express an existing solution, not a promise to construct one.

\section{Uniform statement and the order of quantifiers}
\begin{theorem}[Uniform products with a static one-dimensional factor]
\label{thm:main}
Let $M$ be a compact manifold with the hypotheses above, and let $g(t)$ be a
smooth Ricci flow on a nondegenerate closed interval $[a,b]$. Let $N$ be a
supplied smooth boundaryless Hausdorff sigma-compact one-dimensional manifold.
There exist $B,C\ge0$ with the following property. For every pair of index types
$\Lambda,\mathcal I$, every family of smooth Riemannian metrics
$h_\ell$ on $N$, and every family of smooth immersed closed curve-shortening
flows $c_{\ell,i}$ in $(M\times N,g(t)\oplus h_\ell)$, the following holds.
For every $L_0,\Theta_0\in\R$, if
\[
 \forall\ell,i,\qquad L_{\ell,i}(a)\le L_0,
 \qquad\Theta_{\ell,i}(a)\le\Theta_0,
\]
then, for all $\ell,i$ and all $t\in[a,b]$,
\begin{equation}\label{eq:main}
 L_{\ell,i}(t)\le L_0e^{B(t-a)},\qquad
 \Theta_{\ell,i}(t)\le(\Theta_0+L_0)e^{C(t-a)}.
\end{equation}
Each $h_\ell$ is independent of Ricci-flow time. No regularity in either index,
and no compactness of $N$, is required.
\end{theorem}
\leanref{CurveControl.Geometry.exists_uniform_static_product_family_estimate}{CurveControl/Geometry/UniformProductEstimates.lean}

The constants are selected before $\Lambda,\mathcal I,h,c,L_0,\Theta_0$. The
proof obtains them from the compact base metric family and the time interval;
it does not inspect the factor metric or the curves when choosing witnesses.
The fixed factor manifold and model are parameters of the Lean context. The
witness construction itself uses only the base. Initial bounds remain explicit
inputs; they are not extracted from a compact parameter family by this module.

\begin{corollary}[Positive-scale families]\label{cor:scale}
Theorem~\ref{thm:main} holds with $\Lambda=(0,\infty)$ and arbitrary supplied
static metrics $h_\lambda$. The coefficients do not require a positive lower
bound on $\lambda$, nor continuity or differentiability of $h_\lambda$ in
$\lambda$.
\end{corollary}
\leanref{CurveControl.Geometry.exists_uniform_positive_scale_product_estimate}{CurveControl/Geometry/UniformProductEstimates.lean}

The supplied one-dimensional manifold may in particular be a smooth circle.
The example module constructs the genuine metric $h_\lambda=\lambda^2h_*$ from
any supplied smooth Riemannian metric $h_*$, proves
$|V|_{h_\lambda}=\lambda|V|_{h_*}$, and instantiates the uniform theorem for all
positive scales. A particular \path{AddCircle} atlas and its starting metric
remain caller inputs. The same example gives bounds uniform over the entire
time interval, with right sides $L_0e^{B(b-a)}$ and
$(\Theta_0+L_0)e^{C(b-a)}$.
\leanref{CurveControl.Examples.StaticFactorApplication.scaledFactorMetric}{CurveControl/Examples/StaticFactorApplication.lean}
\leanref{CurveControl.Examples.StaticFactorApplication.exists_scaledFactor_fullInterval_bounds}{CurveControl/Examples/StaticFactorApplication.lean}
The proof of Theorem~\ref{thm:main} is the combination of the following sections.

\section{The corrected local evolution}
The correction of Morgan and Tian identifies the omitted time variation of the
ambient connection. It produces an error linear in $k$, as well as quadratic
terms, and consequently a total-curvature estimate involving initial length
\cite{morgantian2015}. The implementation retains this contribution but uses a
spatial formulation instead of constructing the full space-time metric.

Fix a valid spatial chart near the point of computation. For a vector field
along $c$, use
\[
 D_tW=\partial_tW+\Gamma_{g(t)}(H,W),\qquad
 D_xW=\partial_xW+\Gamma_{g(t)}(X,W).
\]
In particular $D_t$ is not the metric-compatible derivative of $g(t)+dt^2$.
Its product rule contains the ambient metric variation:
\begin{equation}\label{eq:moving}
 \partial_t\ip VW=-2\Ric(V,W)+\ip{D_tV}W+\ip V{D_tW}.
\end{equation}
Let $A=\partial_t\nabla^{g(t)}$, with the spatial coordinate held fixed. The
actual connection variation theorem gives
\begin{equation}\label{eq:A}
 \ip{A(U,V)}W=-(\nabla_U\Ric)(V,W)
 -(\nabla_V\Ric)(U,W)+(\nabla_W\Ric)(U,V).
\end{equation}
This is the lowering of $\partial_t\Gamma$, not the derivative of a lowered
Christoffel symbol. Equation~\eqref{eq:moving} supplies the separate derivative
of the metric.

For signs, let $R^+(U,V)W=D_UD_VW-D_VD_UW-D_{[U,V]}W$.
The upstream do Carmo convention is $R^-=-R^+$. In the evolution formula write
\[
 \operatorname{rm}(U,V,W,Z)=\ip{R^-(U,V)W}Z.
\]
The last-slot antisymmetry yields
$\operatorname{rm}(H,S,H,S)=\ip{R^+(H,S)S}H$.
The coordinate-to-intrinsic bridge proves this equality with its signs; it does
not identify the two curvature operators.

Put $\alpha=q+\Ric(S,S)$ and $P=D_sH+qS$. The velocity and connection
commutators give
\[
 v_t=-\alpha v,\quad D_tS=D_sH+\alpha S,\quad
 D_tH=D_s^2H+2\alpha H+(D_s\alpha)S+R^+(H,S)S+A(S,S).
\]
Since $\ip HS=0$ and $\ip{D_sH}S=-q$, $P$ is the spatial normal component of
$D_sH$, and $|D_sH|^2=|P|^2+q^2$.

\begin{proposition}[Actual squared-curvature evolution]\label{prop:q}
For an existing smooth immersed closed curve-shortening flow in a smooth
ambient Ricci flow on a time interval $J$, at each $t\in\operatorname{interior}J$
and each parameter $x$, one has
\begin{align}\label{eq:q}
 q_t={}&D_s^2q-2|P|^2+2q^2+4q\Ric(S,S)-2\Ric(H,H)\notag\\
 &+2\operatorname{rm}(H,S,H,S)
 -4(\nabla_S\Ric)(S,H)+2(\nabla_H\Ric)(S,S).
\end{align}
\end{proposition}
\leanref{CurveControl.Geometry.IsCurveShorteningFlowOn.hasDerivAt_curvatureSq_time}{CurveControl/Geometry/CurvatureEvolution.lean}

The hypotheses of this declaration include the actual two flows, the geometric
context in Appendix~\ref{app:interfaces}, an interior time, and a parameter.
They do not include a proposed curvature PDE, a Kato inequality, a connection
variation identity, or differentiability of $H$ as additional assumptions.
The latter facts are derived using fixed-chart neighborhoods and equality of
function germs. Choosing the chart centered at the current point occurs only
after the differentiated identities have been established.

The relation to the space-time projection can be checked explicitly. If
$T=\partial_t$ and $\beta=\Ric(S,H)$, that projection is
$\widehat P=P+\beta T$. Its squared norm contains $\beta^2$, absent from the
spatial $P$. The corresponding space-time curvature terms cancel this extra
square and the mixed quadratic Ricci term; the remaining connection derivatives
are precisely the last two terms of \eqref{eq:q}. This explains why copying a
space-time projection into a spatial formula would be incorrect.

\section{Regularization and the moving integral}
Let $B,K,D\ge0$ bound the ambient Ricci tensor, curvature tensor and covariant
derivative of Ricci, respectively, in the evolving metric.
A bound is an estimate on evaluation on arbitrary vectors, by the product
of their metric norms; it is not a bound on unrelated chart coordinates.
Because $|S|=1$ and $|H|=k$, the last five terms of \eqref{eq:q} have absolute
value at most
\[
 (6B+2K)q+6D\sqrt q\le A_0(q+\sqrt q),\qquad
 A_0=\max(6B+2K,6D).
\]
The geometric Kato estimate is $(D_sq)^2\le4q|P|^2$.
For $\varepsilon>0$, define $h_\varepsilon=\sqrt{q+\varepsilon^2}$.
Differentiation, the Kato estimate, and the nonpositive dissipation term give
\begin{equation}\label{eq:reg}
 (h_\varepsilon)_t\le D_s^2h_\varepsilon+k^3+
 \frac{A_0}{2}(h_\varepsilon+1).
\end{equation}
There is no differentiation of $k$ at a zero of $q$.
\leanref{CurveControl.Analysis.deriv_regularizedNorm_le}{CurveControl/Analysis/RegularizedPDE.lean}

Let $r=\Ric(S,S)$ and
$\Theta_\varepsilon=\int_0^1 h_\varepsilon v\,dx$.
The true density equation is $v_t=-(q+r)v$. Inserting it in the differentiated
integral, using periodicity to eliminate the diffusion flux, and using
$k^3-h_\varepsilon q\le0$ yields
\[
 L'\le BL-E,\qquad
 \Theta_\varepsilon'\le(A_0/2+B)\Theta_\varepsilon+(A_0/2)L.
\]
Thus $C=A_0/2+B$ controls $\Theta_\varepsilon+L$. Parameter differentiation,
continuity, integrability, and the periodic endpoint equality are constructed
from the actual smooth functions. The geometric caller does not supply them.
\leanref{CurveControl.Analysis.ScalarCurveHypotheses.regularized_comparison}{CurveControl/Analysis/CurveEstimateAssembly.lean}

The comparison argument uses derivatives only at interior times and continuity
on the closed time interval. In particular it does not claim a two-sided time
derivative at the initial endpoint. For each fixed time,
\[
 0\le\Theta_\varepsilon-\Theta\le\varepsilon L.
\]
Taking limits in the already established comparison, at both its initial and
final times, avoids exchanging a limit with a time derivative.
\leanref{CurveControl.regularizedWeightedCurvature_sub_bounds}{CurveControl/Analysis/RegularizedIntegral.lean}

\begin{theorem}[Geometric estimate with explicit ambient bounds]\label{thm:base}
Let $g,c$ be the actual flows described above on $J$. Let $a\le b$ and
$[a,b]\subseteq J$, and suppose the actual nonnegative ambient tensor bounds $B,K,D$ hold
at every point and time of $J$. With $C=\max(6B+2K,6D)/2+B$, for all
$t\in[a,b]$,
\[
 L(t)\le L(a)e^{B(t-a)},\qquad
 \Theta(t)+L(t)\le(\Theta(a)+L(a))e^{C(t-a)}.
\]
\end{theorem}
\leanref{CurveControl.Geometry.IsCurveShorteningFlowOn.curve_estimate}{CurveControl/Geometry/CurveEstimates.lean}

The intermediate structure \path{CurveControl.Analysis.ScalarCurveHypotheses} is an
explicit conditional analysis interface containing a squared-curvature PDE,
Kato, and density assumptions. The geometric theorem proves its fields from
the same $g,c$; passing that structure alone is not the geometric result.

\begin{proposition}[Accumulated curvature energy]\label{prop:energy}
Under the same flow and interval hypotheses, it is enough to assume
$\Ric(V,V)\ge-B$ for every unit vector ($g(t)(V,V)=1$), at every point and time of $[a,b]$,
to obtain
\[
 \int_a^b e^{-B(t-a)}E(t)\,dt
 \le L(a)-e^{-B(b-a)}L(b).
\]
If also $B\ge0$, then
$\int_a^b E(t)\,dt\le e^{B(b-a)}L(a)-L(b)$.
\end{proposition}
\leanref{CurveControl.Geometry.IsCurveShorteningFlowOn.weighted_curvatureEnergy_integral_le}{CurveControl/Geometry/LengthEstimates.lean}
\leanref{CurveControl.Geometry.IsCurveShorteningFlowOn.curvatureEnergy_integral_le}{CurveControl/Geometry/LengthEstimates.lean}

The blueprint's integral-evolution node additionally states local absolute
continuity of unregularized $\Theta$ and an almost-everywhere differential
inequality. Those statements have not been implemented literally here. The
regularized comparison followed by a limit is a sufficient alternative for
Theorem~\ref{thm:base}. A downstream argument requiring that AC/a.e. interface
would need a further theorem.

\section{Ambient constants from compactness}
\begin{proposition}[Finite actual tensor bounds]\label{prop:compact}
Let $M$ be compact with the geometric context of
Appendix~\ref{app:interfaces}, and let $g$ be a jointly smooth metric family on
a time set $J$. For $[a,b]\subseteq J$ there exist $B,K,D\ge0$ bounding the
actual $\Ric$, $\operatorname{rm}$ and $\nabla\Ric$ everywhere on
$M\times[a,b]$, in the multilinear metric-norm sense used above.
No Ricci-flow equation is needed for this proposition.
\end{proposition}
\leanref{CurveControl.Geometry.CompactAmbientBounds.exists_uniformTensorBoundsOn}{CurveControl/Geometry/CompactAmbientBounds.lean}

The proof passes through finitely many compact chart pieces. Joint smoothness
provides continuity of the relevant tensor coefficients and the inverse Gram
matrix. On each compact piece, finite coefficient families are bounded. Metric
norm control of the coordinates converts these coefficient bounds into intrinsic
multilinear estimates. A finite cover combines the local nonnegative constants
into global ones. This step avoids assuming an unproved globally continuous
operator-norm bundle section. All coefficients in the construction are those of
the canonical geometric tensors; no surrogate curvature field is supplied.

Applying this proposition to the base Ricci flow fixes $B,K,D$ independently
of every later auxiliary metric. What remains is to preserve the same constants
when passing to a product, without dependence on the geometry of the small
factor. Compactness of the product, or a uniform lower injectivity-radius bound,
is not used for this purpose.

\section{The genuine product and one-dimensional uniformity}
The product construction starts with the upstream metric
\path{DCProductMetric}. A technical distinction matters: the ordinary norm on
$E\times F$ is a maximum norm, whereas the canonical intrinsic interfaces
require an inner-product model. The development uses the continuous linear
equivalence to the L$^2$ type \path{WithLp 2} on $E\times F$ and transports the model with corners.
It constructs the induced genuine metric on the same product manifold.

This is not an assertion that the maximum norm is Hilbertian. Extended charts,
tangent-coordinate changes, and the chart Gram pairing are transported
explicitly. Differentiating the Gram identity and applying the algebraic
uniqueness of a symmetric metric-compatible connection proves the actual
Christoffel formula. Differentiating that formula on an open chart target
then proves the curvature formula.
\leanref{CurveControl.Geometry.chartChristoffelContraction_productL2}{CurveControl/Geometry/ProductL2Connection.lean}
\leanref{CurveControl.Geometry.chartCurvature_productL2}{CurveControl/Geometry/ProductL2Connection.lean}

For intrinsic curvature, an isometry is built from the actual tangent-fiber
metrics, not from model-space norms. The four-tensor splits into the sum of
the two factor tensors. Ricci splitting follows by tracing in orthonormal
bases under this actual isometry. Finally the covariant Ricci derivative is
identified with $\partial\Ric-\Gamma\Ric-\Gamma\Ric$ in coordinates;
the derivative of the lower-order splitting is justified on a neighborhood.
\leanref{CurveControl.Geometry.canonicalCurvatureFormAt_productL2}{CurveControl/Geometry/ProductIntrinsicCurvature.lean}
\leanref{CurveControl.Geometry.ricciTensorAt_productL2}{CurveControl/Geometry/ProductIntrinsicCurvature.lean}
\leanref{CurveControl.Geometry.ProductCovRicci.covRic_productL2}{CurveControl/Geometry/ProductCovRicci.lean}

\begin{proposition}[Preservation of the base constants]\label{prop:product}
Let $g$ and $h$ be genuine smooth metrics on the manifolds $M,N$ above, with
$\dim N=1$. If $B,K,D$ bound the three actual base tensors at $p\in M$,
the same constants bound the product tensors at every $(p,z)$.
For a base Ricci flow and static $h$, $g(t)\oplus h$ is itself a Ricci flow
on the same time interval.
\end{proposition}
\leanref{CurveControl.Geometry.productL2_tensorBounds_of_finrank_one}{CurveControl/Geometry/UniformProductEstimates.lean}
\leanref{CurveControl.Geometry.isRicciFlowOn_productL2_oneDimensional}{CurveControl/Geometry/ProductFlow.lean}

Indeed, antisymmetry of the algebraic curvature tensor and collinearity of all
vectors in dimension one give zero canonical curvature. Consequently the
actual Ricci field vanishes identically; differentiating this field gives zero
covariant Ricci. It would not suffice to infer a derivative from vanishing at a
single point. The second-factor terms in all product splittings therefore
vanish. Each projection decreases the actual metric norm, so the base constants
are unchanged.
\leanref{CurveControl.Geometry.canonicalCurvatureFormAt_eq_zero_of_finrank_one}{CurveControl/Geometry/OneDimensionalCurvature.lean}
\leanref{CurveControl.Geometry.covRicciAt_eq_zero_of_finrank_one}{CurveControl/Geometry/OneDimensionalCurvature.lean}

For the time-dependent statement, joint within-smoothness of the product metric
is proved by pulling back the base family and the static factor form along the
two projections and adding them. The within-time derivative has no factor
contribution. Actual Ricci splitting identifies the resulting derivative with
$-2\Ric$ of the product. Thus product flow status is a conclusion.
Combining Propositions~\ref{prop:compact} and~\ref{prop:product} with
Theorem~\ref{thm:base} proves Theorem~\ref{thm:main}; the initial bound for
$\Theta+L$ is $\Theta_0+L_0$, and $L(t)\ge0$ yields \eqref{eq:main}.

\section{Implementation, provenance, and limits}
The dependency base is mathlib at commit
\path{520045ab14e26149ee970e2e617ca04b09bde5d6} and the frenzymath source tree
at commit \path{bb91a091f0b968f8bbe8d861e025a88d82b161be}
\cite{mathlib,frenzymath}. The project directly imports MorganTianLib, which
in turn imports DoCarmoLib. These libraries provide the smooth Riemannian
metric, canonical Levi--Civita connection, chart tensor, and Ricci-flow
infrastructure. mathlib provides the general manifold, calculus, linear
algebra, and integration framework. The newly supplied development consists
of the curve-specific definitions and bridges, spatial evolution proof,
regularized integral chain, compact tensor bounds, product transport and
splitting, and final quantified integration of these parts. We make no claim
of a new mathematical estimate or of priority for its formalization.

The frozen blueprint's labels
\path{thm:curve-flow-curvature-evolution},
\path{lem:regularized-curve-curvature-inequality}, and
\path{thm:curve-flow-length-total-curvature-bounds} describe the corresponding
mathematical goals. The time-space connection and AC/a.e. intermediate nodes
are handled by the alternatives explained above; a matching final conclusion
should not be confused with literal implementation of every blueprint node.

The release record reports a successful complete target build with Lean 4.32.1
and an environment-level reachable-axiom check of the project declarations.
The reachable axioms were only \path{propext}, \path{Classical.choice}, and
\path{Quot.sound}; no \path{sorryAx} or additional axiom was found. Dependency
fingerprints and source hashes are bound to the release record. The build reused
dependency caches, so it is not a claim of reproduction on a clean machine.
The machine check establishes correctness relative to the definitions and
axioms; matching the mathematical specification also requires semantic review.

AI assistance was used in development. Separate Astra agents reviewed proofs
they had not authored; these reviews were not external human peer review.

The principal limitations are explicit inputs and omitted neighboring modules:
existence and continuation of the smooth Ricci and curve flows; construction
of concrete circle atlases and their starting metrics; construction of ramps
and their uniform initial bounds; and the remaining geometric and topological
steps of finite extinction. No acceptance by an external mathematical team is
asserted. The delivered interface can be used once those objects and hypotheses
have been supplied.

\appendix
\section{Exact interface assumptions}\label{app:interfaces}
All declaration names in this report are fully qualified. The main product
roots use the following context; symbols are printed in ASCII to keep the
inventory independent of a Lean font setup.
\begin{quote}\small\ttfamily
E, F : Type*; H, K, M, N : Type*\\
NormedAddCommGroup E; InnerProductSpace Real E\\
NormedAddCommGroup F; InnerProductSpace Real F\\
FiniteDimensional Real E; FiniteDimensional Real F\\
NeZero (Module.finrank Real E)\\
NeZero (Module.finrank Real F)\\
TopologicalSpace H; TopologicalSpace K\\
I : ModelWithCorners Real E H\\
J : ModelWithCorners Real F K\\
I.Boundaryless; J.Boundaryless\\
TopologicalSpace M; ChartedSpace H M; IsManifold I infinity M\\
TopologicalSpace N; ChartedSpace K N; IsManifold J infinity N\\
SigmaCompactSpace M; T2Space M\\
SigmaCompactSpace N; T2Space N; CompactSpace M
\end{quote}
The explicit arguments are a metric family $g$ of type
\path{Real -> Riemannian.RiemannianMetric I M},
$a,b:\R$, \path{MorganTianLib.IsRicciFlowOn g (Set.Icc a b)}, $a\le b$, and
\path{Module.finrank Real F = 1}. In the inventory, \texttt{Real} and
\texttt{infinity} stand for Lean's real type and smoothness order $\infty$.
The base estimate uses only the $E,H,I,M$ part of the context, without
\path{CompactSpace M}, and takes actual uniform tensor bounds as input.
The compact-bound theorem also lists \path{CompleteSpace E}; this instance
is available from finite dimensionality over the complete real field, and is
inferred when the final product root calls it.

\path{MorganTianLib.IsRicciFlowOn} has four fields: order-connectedness and
nontriviality of the time set, joint smoothness of the metric as a horizontal
bilinear-form section, and the within-time equation evaluated on every point
and pair of tangent vectors. Hence its specialization to $[a,b]$ rules out a
singleton interval even though the root also lists $a\le b$.

\path{CurveControl.Geometry.IsCurveShorteningFlowOn} requires joint
\path{ContMDiffOn} smoothness on the time set times $\R$, period one at each
time in that set, nonzero actual velocity at every such time and parameter,
and a \path{HasDerivWithinAt} equation for the local time curve whose derivative
is the actual curvature vector. These are the complete substantive flow
conditions; no target inequality is a field.

The final quantifier order, suppressing only this fixed context, is
\[
 \begin{gathered}
 \exists B,C\ge0\quad\forall\Lambda,\mathcal I,h,c,\\
 \bigl(\text{actual curve flows}\bigr)\ \Longrightarrow
 \forall L_0,\Theta_0,\\
 \bigl(\text{initial bounds}\bigr)\ \Longrightarrow
 \forall\ell\in\Lambda\;\forall i\in\mathcal I\;\forall t\in[a,b]\quad\eqref{eq:main}.
 \end{gathered}
\]
The types $\Lambda,\mathcal I$ are arbitrary; an empty family is allowed.
The positive-scale root fixes $\Lambda$ to the subtype
$\{\lambda:\R\mid0<\lambda\}$ before quantifying over the supplied metric
family. The factor metric has no time argument. Neither a curvature-splitting
hypothesis nor a product Ricci-flow hypothesis appears in either final root.

\bibliographystyle{plain}
\bibliography{references}
\end{document}
